\section{Introduction}

Feature normalization is part of the learning procedure. Dividing a predictor by
its standard deviation changes how quickly gradient methods learn its
coefficient and how strongly individual observations perturb that coefficient.
Under a finite training budget, these changes can affect the predictor even when
the unregularized least-squares solution is invariant to scaling. The
statistical question is therefore how normalization changes prediction at the
available budget, and which effects remain after calibrating the other training
choices.

Rare features make this question concrete. Without scaling, a rare binary
feature has small curvature and can remain poorly learned when common features
are already fitted. Standardization increases its learning rate relative to
common features. It also magnifies the stochastic contribution of the few
observations on which the feature is present. Early stopping and batch size
determine the balance between those effects. Dependence adds a further issue:
prediction for a rare subgroup depends on joint coefficient errors, so the
ordering of individual coefficient errors need not determine the ordering of
subgroup prediction risks.

We analyze squared-loss linear regression with a common and a rare binary
predictor. The population comparisons use known centering and scales, zero
initialization, a fixed intercept, and fresh observations at each update. Our
main result compares the lowest rare-presence risk attainable by each
normalization over all nonnegative constant learning rates and deterministic
stops within a common observation budget. Standardization wins on a specified
rare-signal family, and unscaled training wins on a common-signal family. Both
gaps are uniform as prevalence decreases with the expected number of rare
training observations held of order one. Thus relative scaling can change
attainable prediction risk even after separate scalar retuning. The theorem
concerns two independent features and one observation per update; it gives no
ordering over learning-rate schedules or data-dependent selection rules.

Supporting comparisons explain particular training paths. Standardization has
lower rare-presence risk at an early stop and higher risk later under a specified
rate calibration; a larger batch changes that ordering. Dependence can also make
aggregate and rare-presence risk favor opposite normalizations, even when
scaling increases both individual coefficient mean squared errors.

The existing experiments motivated the main theorem and examine the limits of
the supporting comparisons. A broader search over rates and stopping removes
both original anchor comparisons, although other laws retain reversals or
objective disagreements. Sampled-data experiments then
replace population preprocessing by estimated centering and scales, fit an
intercept, and select configurations on validation data. A sparse categorical
task shows a rare-subgroup benefit within the declared search, alongside
rate-boundary selections and a tradeoff with aggregate prediction. Missing rare
validation observations can conceal extreme subgroup error. These findings
motivate evaluating normalization as part of the full training and selection
procedure. They do not constitute a held-out test of the subsequently developed
oracle prediction.

\paragraph{Relation to prior work.}

The connection between preconditioning and statistical performance is
established \citep{amari2020preconditioning, su2024preconditioning}.
\citet[Section 4.2, Proposition 7]{amari2020preconditioning} relate preferred
preconditioning to signal alignment even after optimizing early-stopping bias
in an overparameterized fixed-data model. Our comparison concerns total
rare-subgroup risk, including stochastic-update variance, in a fixed-dimensional
fresh-sample model with an equal observation cap. It separates both complete
constant-rate and stopping families by explicit gaps. Early
stopping changes least-squares regularization \citep{ali2019early};
least-mean-squares analyses already separate bias and variance and study
sampling and batching \citep{defossez2015averaged, jain2018parallelizing}.
AdaBatch relates sparse mini-batch aggregation to coordinate rescaling
\citep{defossez2017adabatch}, and exact mean-square stability can be expressed
through moment operators \citep{mulayoff2024exact}. Target-aware preconditioning
under covariate shift also studies the dependence on the test objective
\citep{liu2025precondition}. For explicit regularization,
\citet{larsson2025normalization} analyze the influence of feature scaling on
shrinkage. Our comparisons specialize these broader mechanisms to conventional
normalization, final iterates, original-scale coefficients, and rare-presence
prediction. Neither diagonal preconditioning nor the existence of a
bias--variance tradeoff is a new mechanism here. The question is what
quantitative comparisons they yield, and how those comparisons change when the
entire training procedure is tuned.
